% вариант сделан специально для тренировочного экзамена 2026.04.04, автор: Горелова

\begin{problem}{1}
Найдите значение выражения:
\end{problem}

\samerowans{\begin{problem}{1а}
$\dfrac{3^{8} \cdot 4^{5} + 6^{8} \cdot 60}{9^{4} \cdot 2^{12} - 6^{9}}$
\end{problem}}

\smallskip

\samerowans{\begin{problem}{1б}
$\dfrac{16^{k+1} - 2^{k+4}}{2^{k+2} \cdot (8^k - 1)}$
\end{problem}}

\smallskip

\samerowans{\begin{problem}{2}
На координатном луче отмечено несколько точек, координаты которых являются натуральными числами. Известно также, что сумма этих чисел равна~64. Если мы каждую точку переместим на четыре единичных отрезка, то сумма координат новых точек будет равняться~96. Сколько точек было отмечено на координатном луче?
\end{problem}}

\smallskip

\begin{problem}{3}
Разложите на неразложимые множители (чтобы дальнейшее разложение было невозможно):
\end{problem}

\bigskip

\noindent \textbf{3а.} $5b + a(b - 3) - 15 + 9b(3 - b)$

\smallskip

\rowans{1ex}

\bigskip

\noindent \textbf{3б.} $-18b^{2}q + 21q^{3} + 12b^{4} - 14b^{2}q^{2}$

\smallskip

\rowans{1ex}

\bigskip

\noindent \textbf{3в*.} $5x^{2} + 3x - 2$

\smallskip

\rowans{1ex}

\smallskip

\begin{problem}{4}
Упростите выражение и найдите его значение:
$$(-2{,}5nk)^{3} \cdot \left(-\dfrac{2}{5}n^{2}k\right)^{2} \quad \text{при } n = -1; \quad k = 0{,}4.$$
\end{problem}

\rowans{1ex}

\bigskip

\samerowans{\begin{problem}{5а}
Решите уравнение:
$$27x^{3} - (4 - 2x)(2x + 4) + 7x = (1 + 3x)(9x^{2} + 1 - 3x) - (1 - 4x)(x + 3)$$
\end{problem}}

\bigskip

\begin{problem}{5б}
Решите уравнение:
$$3|3x - 2| - 7 = 14 - |2 - 3x|$$
\end{problem}


\noindent\grid{34}{12}

\smallskip


\begin{problem}{6}
Какое из чисел $A$ и $B$ больше и на сколько:
$$A = \underbrace{888\ldots88}_{31\text{ цифра}} \cdot \underbrace{333\ldots33}_{43\text{ цифры}}
\text{ или } B = \underbrace{444\ldots44}_{31\text{ цифра}} \cdot \underbrace{666\ldots67}_{43\text{ цифры}}$$
\end{problem}

\smallskip

\rowans{1ex}

\smallskip

\begin{problem}{7}
Дана линейная функция $y = 3x + b$. При каком значении~$b$ график этой функции:
\end{problem}

\smallskip

\samerowans{\begin{problem}{7а}
пересекает ось $Ox$ в точке с абсциссой~$2$?
\end{problem}}

\smallskip

\samerowans{\begin{problem}{7б}
проходит через точку $P(-3;\; 2)$?
\end{problem}}

\smallskip

\samerowans{\begin{problem}{7в}
пересекает ось ординат в точке с ординатой~$-4$?
\end{problem}}

\smallskip

\samerowans{\begin{problem}{7г}
проходит через точку пересечения графиков функций $y = 0{,}5x + 3$ и $y = 2x - 3$?
\end{problem}}

\smallskip

\begin{problem}{8}
Напишите уравнение линейной функции $y = kx + b$, график которой параллелен графику функции $y = x^6 - (x^3 - 1)(x^3 + 1) - 0{,}25(8 - 12x)$ и проходит через начало координат. Постройте этот график.
\end{problem}

\smallskip

\noindent\grid{34}{24}

\smallskip

\begin{problem}{9}
Два велосипедиста выехали из посёлка в одном направлении: первый~--- в~9~часов утра, а второй~--- в~11~ч.~30~мин. После того как второй велосипедист догнал первого, они ещё продолжили путь в течение 3~часов. В~момент остановки оказалось, что второй велосипедист обогнал первого на~24~км. В~котором часу второй велосипедист догнал первого и какое расстояние проехали велосипедисты до этого, если скорость первого велосипедиста равна~$0{,}6$ скорости второго?
\end{problem}

\smallskip

\noindent\grid{34}{18}

\smallskip

\begin{problem}{10}
    На конкурсе выступление каждого участника оценивают семь судей. Каждый судья выставляет оценку~--- целое число баллов от~0 до~10 включительно. Известно, что за выступление участника~<<A>> все члены жюри выставили различные оценки. По старой системе оценивания итоговый балл за выступление определяется как среднее арифметическое всех оценок судей. По новой системе оценивания итоговый балл вычисляется следующим образом: отбрасываются наименьшая и наибольшая оценки и подсчитывается среднее арифметическое пяти оставшихся оценок. В~результате оказалось, что разность итоговых баллов, вычисленных по старой и новой системам оценивания, равна~$\dfrac{2}{35}$. Какие оценки могли быть выставлены судьями участнику~<<A>>?
\end{problem}

\bigskip

\rowans{1ex}

\bigskip

\begin{problem}{11 (12 баллов)}
Решите ОДНО из заданий по выбору.

\smallskip

\noindent а) При каких $a$ уравнение $ax-6(1+x)=2x+5-0{,}5a$ имеет корни?

\smallskip

\noindent б) Чему равна сумма $\dfrac{1}{3\cdot7}+\dfrac{1}{7\cdot11}+\dots+\dfrac{1}{95\cdot99}+\dfrac{1}{99\cdot103}$?

\smallskip

\noindent в) Можно ли найти такое число $n$, что $n^3-n$ равно 1514?
\end{problem}

\smallskip

\noindent\grid{34}{24}